\chapter{Dimensi Lanjutan dan SCD Type 2}
\label{modul:dimensi-lanjutan-scd2}

\section{Identitas Modul}

\begin{identitasmodul}
\textbf{Sumber materi}    & Pertemuan 4 \\
\textbf{CPMK}             & CPMK-092 \\
\textbf{Minggu pelaksanaan} & Minggu ke-5 \\
\textbf{Durasi tatap muka} & 120 menit \\
\textbf{Estimasi tugas}   & 3--4 jam di luar sesi \\
\textbf{Moda kerja}       & Individual \\
\textbf{Perangkat}        & PostgreSQL 17 dan ekstensi \perintah{btree\_gist} \\
\textbf{Dataset}          & NusaMart ukuran kecil dengan skenario perubahan produk \\
\textbf{Skrip pendukung}  & \berkas{sql/modul-03/} beserta \berkas{check.sql} \\
\textbf{Luaran}           & DDL SCD-2, skrip pemutakhiran, dan bukti tiga versi satu produk \\
\end{identitasmodul}

\slideref{4}

\section{Capaian dan Indikator Keberhasilan}

Mahasiswa mampu mempertahankan histori dimensi dengan SCD Type 2, mencegah
rentang versi yang tumpang tindih, dan menerapkan role-playing serta junk
dimension. Satu produk harus memiliki histori yang benar tanpa dua baris kini.

\begin{capaian}
Pada akhir modul ini mahasiswa mampu:
\begin{enumerate}[leftmargin=1.4em]
  \item membedakan SCD Type 0, Type 1, dan Type 2, serta menetapkan perlakuan
        yang tepat untuk setiap atribut dimensi;
  \item mengubah \perintah{dim\_produk} menjadi SCD Type 2 dengan kolom
        \perintah{mulai\_berlaku}, \perintah{selesai\_berlaku}, dan
        \perintah{baris\_kini};
  \item memasang \istilah{exclusion constraint} sehingga basis data sendiri
        menolak versi yang tumpang tindih;
  \item memuat perubahan dengan menutup versi lama dan membuka versi baru,
        bukan menimpa;
  \item menulis query yang mengembalikan versi berlaku pada tanggal tertentu,
        bukan versi terakhir;
  \item merancang \istilah{junk dimension} dan memakai satu dimensi dalam
        beberapa peran.
\end{enumerate}
\end{capaian}

Modul ini memperbaiki satu cacat yang sengaja dibiarkan pada Modul 2. Dimensi
yang Anda bangun minggu lalu hanya menyimpan keadaan terkini: begitu harga
sebuah produk berubah, versi lamanya hilang dan laporan tahun lalu ikut
berubah. Gudang data yang angkanya berubah sendiri tanpa ada transaksi baru
tidak dapat dipakai untuk apa pun.

\section{Prasyarat dan Persiapan}

\subsection{Prasyarat}

\begin{itemize}
  \item Skema bintang hasil Modul 2 yang sudah terisi dan terekonsiliasi;
  \item pemahaman pemisahan surrogate key dan natural key, khususnya alasan
        ketiga yang dibahas pada Modul 2;
  \item operasi rentang tanggal: perbandingan \perintah{>=} dan \perintah{<}
        pada \perintah{DATE}, serta pengertian batas terbuka dan tertutup;
  \item \perintah{UPDATE} dan \perintah{INSERT} dengan subquery, serta
        pengertian transaksi basis data.
\end{itemize}

\subsection{Pre-lab}

\begin{enumerate}[leftmargin=1.4em]
  \item Aktifkan ekstensi \perintah{btree\_gist} pada \perintah{nusamart\_dw}.
        Ekstensi ini tersedia pada image Docker praktikum, tetapi harus
        diaktifkan per basis data.
\begin{lstlisting}[style=sqlgaya]
CREATE EXTENSION IF NOT EXISTS btree_gist;
SELECT extname, extversion FROM pg_extension WHERE extname = 'btree_gist';
\end{lstlisting}
  \item Bila Modul 2 belum selesai, pulihkan \berkas{snapshot-modul-02.sql}
        sehingga keempat tabel skema bintang tersedia dan terisi.
  \item Siapkan jawaban atas pertanyaan berikut pada
        \berkas{modul-03-scd2/pre-lab.md}:
  \begin{enumerate}[label=\alph*.,leftmargin=1.4em]
    \item Pada Modul 2, \perintah{produk\_id} dideklarasikan \perintah{UNIQUE}
          di \perintah{dim\_produk}. Setelah dimensi menyimpan beberapa versi
          per produk, apa yang harus terjadi pada constraint itu?
    \item Bila kolom \perintah{nama\_produk} salah eja dan diperbaiki, apakah
          layak dibuatkan versi baru? Jelaskan singkat.
    \item Berapa baris \perintah{dim\_produk} Anda saat ini? Catat angkanya
          sebagai pembanding setelah pemuatan SCD-2.
  \end{enumerate}
\end{enumerate}

\section{Konsep Dasar}

\subsection{Slowly Changing Dimension Type 2}

Atribut dimensi berubah, tetapi jarang. Perubahan itu disebut \istilah{slowly
changing}, dan cara menanganinya menentukan apakah laporan lama tetap benar
\citep{kimball2013}. Ada tiga perlakuan yang perlu Anda kuasai.

\begin{center}
\small
\begin{tabular}{@{}p{0.12\textwidth}>{\raggedright\arraybackslash}p{0.34\textwidth}>{\raggedright\arraybackslash}p{0.46\textwidth}@{}}
\toprule
\textbf{Tipe} & \textbf{Perlakuan} & \textbf{Contoh pada NusaMart} \\
\midrule
Type 0 & Nilai tidak pernah berubah &
  \perintah{sku} dan \perintah{tanggal\_buka} toko \\
Type 1 & Nilai lama ditimpa, histori hilang &
  Perbaikan salah eja \perintah{nama\_produk} \\
Type 2 & Baris baru dibuat, versi lama disimpan &
  Perubahan \perintah{harga\_jual} dan \perintah{kategori} \\
\bottomrule
\end{tabular}
\end{center}

Pilihannya bukan selera. Aturannya: gunakan Type 1 bila nilai lama memang
\emph{salah}, dan Type 2 bila nilai lama pernah \emph{benar}. Salah eja nama
produk tidak pernah benar, jadi cukup ditimpa. Harga tahun lalu pernah benar,
jadi harus disimpan.

\begin{konsep}{Tiga kolom penanda SCD Type 2}
\begin{itemize}
  \item \perintah{mulai\_berlaku} --- tanggal versi ini mulai berlaku.
  \item \perintah{selesai\_berlaku} --- tanggal versi ini berhenti berlaku,
        bersifat \textbf{eksklusif}. Versi kini diberi nilai penjaga
        \perintah{9999-12-31}.
  \item \perintah{baris\_kini} --- penanda \perintah{BOOLEAN} untuk versi yang
        sedang berlaku. Secara logika ia berlebihan, tetapi membuat query
        harian jauh lebih ringkas.
\end{itemize}
Satu versi berlaku pada rentang $[\text{mulai}, \text{selesai})$: tanggal mulai
termasuk, tanggal selesai tidak.
\end{konsep}

Batas eksklusif ini penting dan sering keliru. Bila versi lama ditutup pada
tanggal yang sama dengan tanggal mulai versi baru, dan keduanya memakai batas
tertutup, maka pada tanggal itu \emph{dua} versi berlaku sekaligus --- dan
setiap query yang mencari harga pada hari itu akan menggandakan barisnya.

\begin{peringatan}
Akibat penyimpanan beberapa versi, \perintah{produk\_id} \textbf{tidak boleh
lagi} dideklarasikan \perintah{UNIQUE}. Ini perubahan struktural, bukan
kelalaian: constraint yang benar pada Modul 2 justru menjadi salah pada Modul 3.
Yang harus tetap dijamin adalah aturan yang lebih lemah --- satu produk hanya
boleh memiliki \emph{satu} baris kini.
\end{peringatan}

\subsection{Rentang Validitas dan Exclusion Constraint}

Dua aturan harus dijaga sepanjang umur dimensi:

\begin{enumerate}[leftmargin=1.4em]
  \item satu \perintah{produk\_id} hanya memiliki satu baris dengan
        \perintah{baris\_kini = TRUE};
  \item rentang berlaku antarversi satu \perintah{produk\_id} tidak boleh
        tumpang tindih.
\end{enumerate}

Aturan pertama dijaga dengan \istilah{partial unique index}. Aturan kedua tidak
dapat dinyatakan sebagai \perintah{UNIQUE} biasa, karena yang dilarang bukan
nilai yang sama melainkan rentang yang \emph{beririsan}. Untuk itu PostgreSQL
menyediakan \perintah{EXCLUDE}.

\begin{lstlisting}[style=sqlgaya]
-- Aturan 1: satu baris kini per produk.
CREATE UNIQUE INDEX uq_dim_produk_kini
  ON dim_produk (produk_id) WHERE baris_kini;

-- Aturan 2: rentang versi satu produk tidak boleh beririsan.
ALTER TABLE dim_produk ADD CONSTRAINT ex_dim_produk_rentang
  EXCLUDE USING gist (
    produk_id WITH =,
    daterange(mulai_berlaku, selesai_berlaku, '[)') WITH &&
  );
\end{lstlisting}

Baris itu dibaca begini: tolak dua baris yang \emph{produk\_id-nya sama}
(\perintah{WITH =}) \emph{dan} rentang tanggalnya \emph{beririsan}
(\perintah{WITH \&\&}).

\begin{catatan}
Ekstensi \perintah{btree\_gist} diperlukan justru karena bagian pertama.
Indeks GiST secara asali menangani operator rentang seperti \perintah{\&\&},
tetapi tidak menangani \perintah{=} pada tipe skalar seperti \perintah{INTEGER}.
\perintah{btree\_gist} menjembatani keduanya sehingga satu constraint dapat
memadukan kesamaan dan irisan.
\end{catatan}

Nilai constraint ini terletak pada waktunya bekerja. Tanpa \perintah{EXCLUDE},
kesalahan pemuatan baru ketahuan berbulan-bulan kemudian ketika sebuah laporan
menghasilkan angka dua kali lipat. Dengan \perintah{EXCLUDE}, pemuatan yang
keliru gagal pada detik itu juga.

\subsection{Role-playing Dimension}

Satu tabel dimensi kerap dipakai beberapa kali dalam satu konteks, dengan
makna yang berbeda pada setiap pemakaian. Inilah \istilah{role-playing
dimension}.

Tabel \perintah{promosi} pada sumber memiliki tanggal mulai dan tanggal
selesai. Keduanya adalah tanggal, keduanya menunjuk \perintah{dim\_tanggal}
yang sama --- tetapi menyebutnya ``tanggal'' pada laporan akan
membingungkan. Penyelesaiannya bukan menyalin \perintah{dim\_tanggal} menjadi
dua tabel, melainkan membuat dua \emph{view} atasnya.

\begin{lstlisting}[style=sqlgaya]
CREATE VIEW dim_tanggal_mulai AS
  SELECT tanggal_key AS tanggal_mulai_key,
         tanggal     AS tanggal_mulai,
         nama_bulan  AS bulan_mulai,
         tahun       AS tahun_mulai
  FROM   dim_tanggal;
\end{lstlisting}

View tidak menyalin data, sehingga tidak ada biaya penyimpanan dan tidak ada
risiko kedua salinan berbeda isi. Yang diperoleh adalah nama kolom yang jelas
pada setiap peran.

\subsection{Junk Dimension dan Bridge Table}

Tabel fakta kerap memuat sejumlah atribut indikator berkardinalitas rendah ---
metode bayar, status, penanda anggota. Menyimpannya sebagai kolom teks pada
tabel fakta boros; membuatkan satu tabel dimensi untuk masing-masing
menghasilkan dimensi kerdil yang menambah join tanpa memberi manfaat.

\begin{konsep}{Junk dimension}
Beberapa atribut indikator digabung menjadi \emph{satu} tabel dimensi berisi
kombinasi nilainya. Tabel fakta cukup menyimpan satu kunci. Pada NusaMart:
lima metode bayar $\times$ tiga status $\times$ dua jenis pembeli menghasilkan
paling banyak 30 baris --- satu dimensi yang sangat kecil menggantikan tiga
kolom teks pada jutaan baris fakta.
\end{konsep}

\istilah{Bridge table} menangani persoalan yang berbeda: hubungan
banyak-ke-banyak. Satu promosi mencakup banyak produk, dan satu produk dapat
tercakup beberapa promosi. Hubungan semacam ini tidak dapat dinyatakan dengan
satu kunci asing pada tabel fakta, dan memaksakannya akan menggandakan measure.
Rancangan bridge table menjadi tugas luar laboratorium modul ini.

\section{Studi Kasus dan Protokol Praktikum}

Perubahan harga produk NusaMart digunakan untuk menguji histori. Query harus
mengembalikan versi dimensi yang berlaku pada waktu transaksi, bukan versi
terakhir saat query dijalankan.

Skrip \berkas{sql/modul-03/01-skenario-perubahan.sql} menerapkan tiga
perubahan berurutan pada satu produk terpilih. Setelah ketiganya dimuat,
dimensi harus memuat tiga versi produk tersebut --- dua tertutup dan satu
berlaku.

\begin{center}
\small
\begin{tabular}{@{}p{0.16\textwidth}p{0.20\textwidth}>{\raggedright\arraybackslash}p{0.24\textwidth}>{\raggedright\arraybackslash}p{0.28\textwidth}@{}}
\toprule
\textbf{Versi} & \textbf{Berlaku} & \textbf{Perubahan} & \textbf{Perlakuan} \\
\midrule
1 & sampai 2024-03-01 & keadaan awal & --- \\
2 & 2024-03-01 sampai 2024-08-15 & harga naik & Type 2 \\
3 & sejak 2024-08-15 & kategori dipindah & Type 2 \\
\bottomrule
\end{tabular}
\end{center}

\begin{catatanasisten}
Skrip skenario memilih produk berdasarkan \perintah{sku}, bukan berdasarkan
\perintah{produk\_id}, sehingga produk yang sama terpilih di seluruh kelas
meskipun surrogate key setiap mahasiswa berbeda. Pastikan mahasiswa mencatat
\perintah{produk\_id} hasil pemilihan itu --- angka tersebut dipakai berulang
kali pada Bagian C dan D.
\end{catatanasisten}

\section{Alur Praktikum 120 Menit}

\begin{alurwaktu}
0--15 & Demonstrasi perubahan dimensi dan histori. \\
15--95 & Kerja terbimbing: SCD-2, constraint, role-playing, dan junk dimension. \\
95--110 & Checkpoint tiga versi produk dan query historis. \\
110--120 & Penjelasan tugas bridge table. \\
\end{alurwaktu}

\section{Prosedur Praktikum Terbimbing}

\subsection{Bagian A: Mengubah Dimensi Produk menjadi SCD-2}

\langkah{A.1 Menetapkan kebijakan per atribut}{10 menit}

Sebelum menyentuh DDL, lengkapi tabel kebijakan berikut pada
\berkas{kebijakan-scd.md}. Setiap baris harus disertai alasan satu kalimat.
Tanpa dokumen ini, keputusan pemuatan menjadi tebakan.

\begin{center}
\small
\begin{tabular}{@{}p{0.24\textwidth}p{0.14\textwidth}>{\raggedright\arraybackslash}p{0.50\textwidth}@{}}
\toprule
\textbf{Atribut} & \textbf{Tipe SCD} & \textbf{Alasan} \\
\midrule
\perintah{sku}          & Type 0 & Identitas produk, tidak pernah berubah \\
\perintah{nama\_produk} & Type 1 & \ldots \\
\perintah{merek}        & \ldots & \ldots \\
\perintah{kategori}     & \ldots & \ldots \\
\perintah{harga\_jual}  & \ldots & \ldots \\
\perintah{berat\_gram}  & \ldots & \ldots \\
\bottomrule
\end{tabular}
\end{center}

\langkah{A.2 Menambahkan kolom histori}{10 menit}

\begin{lstlisting}[style=sqlgaya]
-- harga_jual belum ada pada Modul 2; ia ditambahkan sekarang karena
-- justru atribut inilah yang akan berubah.
ALTER TABLE dim_produk ADD COLUMN harga_jual      NUMERIC(12,2);
ALTER TABLE dim_produk ADD COLUMN mulai_berlaku   DATE;
ALTER TABLE dim_produk ADD COLUMN selesai_berlaku DATE;
ALTER TABLE dim_produk ADD COLUMN baris_kini      BOOLEAN;
ALTER TABLE dim_produk ADD COLUMN versi           SMALLINT;

-- Baris yang sudah ada menjadi versi pertama, berlaku sejak awal data.
UPDATE dim_produk
SET    mulai_berlaku   = DATE '2020-01-01',
       selesai_berlaku = DATE '9999-12-31',
       baris_kini      = TRUE,
       versi           = 1,
       harga_jual      = (SELECT sp.harga_jual FROM staging_produk sp
                          WHERE sp.produk_id = dim_produk.produk_id);

ALTER TABLE dim_produk ALTER COLUMN mulai_berlaku   SET NOT NULL;
ALTER TABLE dim_produk ALTER COLUMN selesai_berlaku SET NOT NULL;
ALTER TABLE dim_produk ALTER COLUMN baris_kini      SET NOT NULL;
\end{lstlisting}

\begin{catatan}
Urutan di atas disengaja: kolom ditambahkan sebagai \perintah{NULL},
diisi, baru kemudian dijadikan \perintah{NOT NULL}. Menambahkan kolom
\perintah{NOT NULL} tanpa nilai asali pada tabel yang sudah berisi data akan
langsung gagal.
\end{catatan}

\langkah{A.3 Mengganti constraint keunikan}{5 menit}

\begin{lstlisting}[style=sqlgaya]
-- Constraint Modul 2 kini menghalangi versi kedua.
ALTER TABLE dim_produk DROP CONSTRAINT dim_produk_produk_id_key;

-- Aturan yang benar: satu baris kini per produk.
CREATE UNIQUE INDEX uq_dim_produk_kini
  ON dim_produk (produk_id) WHERE baris_kini;
\end{lstlisting}

Uji bahwa indeks parsial benar-benar bekerja sebelum melanjutkan:

\begin{lstlisting}[style=sqlgaya]
-- Harus GAGAL: dua baris kini untuk satu produk.
INSERT INTO dim_produk (produk_id, sku, nama_produk, mulai_berlaku,
                        selesai_berlaku, baris_kini, versi)
SELECT produk_id, sku, nama_produk, DATE '2025-01-01',
       DATE '9999-12-31', TRUE, 99
FROM   dim_produk WHERE baris_kini LIMIT 1;
\end{lstlisting}

\subsection{Bagian B: Menjaga Rentang Versi}

\langkah{B.1 Memasang exclusion constraint}{10 menit}

\begin{lstlisting}[style=sqlgaya]
ALTER TABLE dim_produk ADD CONSTRAINT ex_dim_produk_rentang
  EXCLUDE USING gist (
    produk_id WITH =,
    daterange(mulai_berlaku, selesai_berlaku, '[)') WITH &&
  );
\end{lstlisting}

\begin{peringatan}
Bila perintah ini gagal dengan pesan \emph{conflicting key value violates
exclusion constraint}, dimensi Anda \emph{sudah} memuat rentang yang tumpang
tindih. Ini kabar baik, bukan kabar buruk: constraint menemukan kesalahan yang
selama ini tersembunyi. Cari barisnya lebih dahulu, perbaiki, lalu pasang
kembali constraint-nya --- jangan menghapus constraint agar perintah berhasil.
\end{peringatan}

\langkah{B.2 Membuktikan constraint menolak irisan}{5 menit}

Jalankan perintah berikut. Ia \textbf{harus} gagal, dan salinan pesan galatnya
menjadi bukti pada laporan.

\begin{lstlisting}[style=sqlgaya]
BEGIN;
-- Versi palsu yang beririsan dengan versi kini produk mana pun.
INSERT INTO dim_produk (produk_id, sku, nama_produk, mulai_berlaku,
                        selesai_berlaku, baris_kini, versi)
SELECT produk_id, sku, nama_produk,
       DATE '2024-06-01', DATE '2024-12-31', FALSE, 98
FROM   dim_produk WHERE baris_kini LIMIT 1;
ROLLBACK;
\end{lstlisting}

Simpan pesan galat pada \berkas{bukti-constraint.md}. Kemampuan menunjukkan
bahwa sebuah aturan \emph{ditolak} basis data lebih meyakinkan daripada
pernyataan bahwa aturan itu dipatuhi.

\subsection{Bagian C: Memuat Perubahan Produk}

\langkah{C.1 Menjalankan skenario perubahan}{5 menit}

Jalankan \berkas{sql/modul-03/\allowbreak 01-skenario-perubahan.sql}. Skrip
menyiapkan dua perubahan pada tabel \perintah{staging\_produk\_perubahan},
masing-masing disertai tanggal berlaku dan nomor \perintah{batch}. Catat
\perintah{produk\_id} yang terpilih --- angka itu dipakai berulang kali sampai
akhir modul.

\langkah{C.2 Menutup versi lama dan membuka versi baru}{20 menit}

Inilah inti modul. Pemuatan SCD-2 selalu terdiri atas dua langkah yang harus
berada dalam \emph{satu} transaksi: menutup versi lama, lalu membuka versi baru.

Perintah berikut dijalankan \textbf{satu batch sekali} --- mula-mula dengan
\perintah{:batch := 1}, lalu diulang dengan \perintah{:batch := 2}.

\begin{lstlisting}[style=sqlgaya]
BEGIN;

-- Langkah 1: tutup versi kini yang atributnya berubah.
UPDATE dim_produk d
SET    selesai_berlaku = p.berlaku_sejak,
       baris_kini      = FALSE
FROM   staging_produk_perubahan p
WHERE  d.produk_id  = p.produk_id
  AND  p.batch      = :batch
  AND  d.baris_kini
  AND  (d.harga_jual IS DISTINCT FROM p.harga_jual
     OR d.kategori   IS DISTINCT FROM p.kategori);

-- Langkah 2: buka versi baru mulai tanggal perubahan.
INSERT INTO dim_produk (produk_id, sku, nama_produk, merek, kategori,
                        sub_kategori, satuan, berat_gram, harga_jual,
                        mulai_berlaku, selesai_berlaku, baris_kini, versi)
SELECT p.produk_id, p.sku, p.nama_produk, p.merek, p.kategori,
       p.sub_kategori, p.satuan, p.berat_gram, p.harga_jual,
       p.berlaku_sejak, DATE '9999-12-31', TRUE,
       COALESCE((SELECT MAX(versi) FROM dim_produk d2
                 WHERE d2.produk_id = p.produk_id), 0) + 1
FROM   staging_produk_perubahan p
WHERE  p.batch = :batch
  AND  NOT EXISTS (SELECT 1 FROM dim_produk d
                   WHERE d.produk_id = p.produk_id AND d.baris_kini);

COMMIT;
\end{lstlisting}

\begin{peringatan}
Cobalah sekali menjalankannya \emph{tanpa} syarat \perintah{p.batch = :batch},
sehingga kedua perubahan dimuat sekaligus. Langkah 1 hanya menutup satu versi,
sedangkan Langkah 2 menyisipkan dua baris kini --- dan indeks parsial
\perintah{uq\_dim\_produk\_kini} menolaknya. Inilah alasan pemuatan SCD-2
dikerjakan per batch menurut urutan waktu: satu produk hanya boleh berpindah
versi satu kali dalam satu jalan pemuatan. Catat pesan galatnya, lalu ulangi
dengan cara yang benar.
\end{peringatan}

\begin{catatan}
Perhatikan \perintah{IS DISTINCT FROM}, bukan \perintah{<>}. Operator biasa
menghasilkan \perintah{NULL} ketika salah satu sisinya \perintah{NULL},
sehingga perubahan dari \perintah{NULL} menjadi sebuah nilai --- atau
sebaliknya --- tidak akan terdeteksi. Ini termasuk kekeliruan SCD-2 yang paling
sulit ditemukan, karena pemuatan berjalan tanpa galat dan hanya sebagian
perubahan yang terlewat.
\end{catatan}

\langkah{C.3 Memverifikasi tiga versi}{5 menit}

\begin{lstlisting}[style=sqlgaya]
SELECT produk_key, produk_id, versi, harga_jual, kategori,
       mulai_berlaku, selesai_berlaku, baris_kini
FROM   dim_produk
WHERE  produk_id = <produk_id yang tercatat pada C.1>
ORDER  BY mulai_berlaku;
\end{lstlisting}

Hasilnya harus memperlihatkan tiga baris: \perintah{selesai\_berlaku} versi
sebelumnya sama persis dengan \perintah{mulai\_berlaku} versi sesudahnya, dan
hanya baris terakhir yang \perintah{baris\_kini}-nya bernilai benar.

\langkah{C.4 Mengarahkan fakta ke versi yang benar}{10 menit}

Pencarian kunci saat memuat fakta kini berubah. Pada Modul 2 satu
\perintah{produk\_id} berpadanan dengan tepat satu baris dimensi; sekarang ia
berpadanan dengan beberapa, dan yang benar ditentukan oleh \emph{tanggal
transaksi}.

\begin{lstlisting}[style=sqlgaya]
-- Bentuk yang benar: versi yang berlaku pada saat transaksi terjadi.
JOIN dim_produk dp
  ON  dp.produk_id = s.produk_id
  AND s.waktu_transaksi::DATE >= dp.mulai_berlaku
  AND s.waktu_transaksi::DATE <  dp.selesai_berlaku
\end{lstlisting}

\begin{peringatan}
Bentuk yang keliru --- dan sangat sering dipakai --- adalah
\perintah{ON dp.produk\_id = s.produk\_id AND dp.baris\_kini}. Query itu
berjalan lancar dan menghasilkan jumlah baris yang tepat, tetapi seluruh
transaksi lama menunjuk harga hari ini. Kesalahan semacam ini tidak
memunculkan galat apa pun; ia hanya membuat laporan tahun lalu berubah setiap
kali ada produk yang harganya naik.
\end{peringatan}

Muat ulang \perintah{fakta\_penjualan} dengan pencarian kunci yang benar, lalu
bandingkan totalnya dengan hasil Modul 2. Jumlah barisnya harus sama persis;
yang berubah hanya \perintah{produk\_key} pada sebagian baris.

\subsection{Bagian D: Menambah Dimensi Khusus}

\langkah{D.1 Membangun junk dimension}{10 menit}

\begin{lstlisting}[style=sqlgaya]
CREATE TABLE dim_transaksi_flag (
  flag_key      INTEGER GENERATED ALWAYS AS IDENTITY PRIMARY KEY,
  metode_bayar  VARCHAR(30) NOT NULL,
  status        VARCHAR(20) NOT NULL,
  jenis_pembeli VARCHAR(15) NOT NULL,   -- 'anggota' atau 'umum'
  UNIQUE (metode_bayar, status, jenis_pembeli)
);

-- Hanya kombinasi yang benar-benar muncul pada data yang dibuat.
INSERT INTO dim_transaksi_flag (metode_bayar, status, jenis_pembeli)
SELECT DISTINCT metode_bayar, status,
       CASE WHEN pelanggan_id IS NULL THEN 'umum' ELSE 'anggota' END
FROM   staging_penjualan;
\end{lstlisting}

\begin{catatan}
Dua pendekatan sama-sama sah: membangkitkan seluruh kombinasi kartesian di
muka, atau hanya kombinasi yang benar-benar muncul. Yang pertama membuat
laporan dapat menampilkan kombinasi berjumlah nol; yang kedua menjaga dimensi
tetap kecil. Untuk NusaMart, selisih keduanya hanya belasan baris --- pilih
salah satu, lalu tuliskan alasannya.
\end{catatan}

Perhatikan bahwa kolom \perintah{jenis\_pembeli} lahir dari temuan Modul 1:
\perintah{pelanggan\_id} yang bernilai \perintah{NULL} bukan kerusakan data,
melainkan pembeli tanpa kartu anggota. Di sinilah fakta bisnis itu memperoleh
tempat yang semestinya.

\langkah{D.2 Membuat role-playing dimension}{10 menit}

\begin{lstlisting}[style=sqlgaya]
CREATE TABLE dim_promosi (
  promosi_key         INTEGER GENERATED ALWAYS AS IDENTITY PRIMARY KEY,
  promosi_id          INTEGER     NOT NULL UNIQUE,
  kode_promo          VARCHAR(20) NOT NULL,
  nama_promosi        VARCHAR(120) NOT NULL,
  tipe_promosi        VARCHAR(40),
  tanggal_mulai_key   INTEGER NOT NULL REFERENCES dim_tanggal(tanggal_key),
  tanggal_selesai_key INTEGER NOT NULL REFERENCES dim_tanggal(tanggal_key)
);

CREATE VIEW dim_tanggal_mulai AS
  SELECT tanggal_key AS tanggal_mulai_key, tanggal AS tanggal_mulai,
         nama_bulan  AS bulan_mulai,       tahun   AS tahun_mulai
  FROM   dim_tanggal;

CREATE VIEW dim_tanggal_selesai AS
  SELECT tanggal_key AS tanggal_selesai_key, tanggal AS tanggal_selesai,
         nama_bulan  AS bulan_selesai,       tahun   AS tahun_selesai
  FROM   dim_tanggal;
\end{lstlisting}

Buktikan manfaatnya dengan satu query yang memakai kedua peran sekaligus:

\begin{lstlisting}[style=sqlgaya]
SELECT p.nama_promosi, m.tanggal_mulai, s.tanggal_selesai,
       s.tanggal_selesai - m.tanggal_mulai AS lama_hari
FROM   dim_promosi p
JOIN   dim_tanggal_mulai   m ON m.tanggal_mulai_key   = p.tanggal_mulai_key
JOIN   dim_tanggal_selesai s ON s.tanggal_selesai_key = p.tanggal_selesai_key
ORDER  BY lama_hari DESC
LIMIT  10;
\end{lstlisting}

Tanpa view, kedua join akan memakai nama kolom yang sama dan setiap query harus
menuliskan alias sendiri. Dengan view, penamaan peran ditetapkan satu kali di
tingkat skema.

\begin{luaran}
\berkas{modul-03-scd2/ddl-scd2.sql}, \berkas{load-scd2.sql},
\berkas{kebijakan-scd.md}, \berkas{bukti-constraint.md}, dan
\berkas{bukti-tiga-versi.md} yang memuat keluaran query verifikasi C.3.
\end{luaran}

\begin{checkpoint}
Asisten menjalankan \berkas{sql/modul-03/check.sql} dan memeriksa butir berikut:
\begin{ceklis}
  \item ekstensi \perintah{btree\_gist} aktif;
  \item \perintah{dim\_produk} memiliki kolom \perintah{mulai\_berlaku},
        \perintah{selesai\_berlaku}, dan \perintah{baris\_kini};
  \item constraint \perintah{UNIQUE} lama pada \perintah{produk\_id} sudah
        dihapus, digantikan indeks parsial;
  \item exclusion constraint terpasang dan menolak rentang yang beririsan;
  \item tidak ada produk dengan lebih dari satu baris kini;
  \item tidak ada celah maupun irisan antarversi satu produk;
  \item produk skenario memiliki tepat tiga versi;
  \item query harga pada tanggal tertentu mengembalikan versi yang benar;
  \item \perintah{dim\_transaksi\_flag} dan \perintah{dim\_promosi} ada beserta
        kedua view perannya.
\end{ceklis}
\end{checkpoint}

\section{Latihan Mandiri di Kelas}

\latihan{Harga pada tiga titik waktu}

Untuk produk skenario, jalankan query harga pada tiga tanggal: sebelum
perubahan pertama, di antara kedua perubahan, dan sesudah perubahan kedua.

\begin{lstlisting}[style=sqlgaya]
SELECT harga_jual, kategori, versi
FROM   dim_produk
WHERE  produk_id = <produk_id skenario>
  AND  DATE '2024-05-20' >= mulai_berlaku
  AND  DATE '2024-05-20' <  selesai_berlaku;
\end{lstlisting}

Lalu jawab:

\begin{enumerate}[leftmargin=1.4em]
  \item Ketiga query mengembalikan berapa baris? Mengapa jumlah itu yang benar?
  \item Ganti klausa rentang dengan \perintah{AND baris\_kini}. Hasil mana yang
        berubah, dan mengapa perubahan itu berbahaya justru karena tidak
        menimbulkan galat?
  \item Apa yang terjadi bila tanggal yang ditanyakan tepat sama dengan
        \perintah{selesai\_berlaku} sebuah versi? Kaitkan dengan batas
        eksklusif.
\end{enumerate}

\latihan{Menghitung ukuran junk dimension}

Hitung berapa baris yang dihemat oleh junk dimension.

\begin{enumerate}[leftmargin=1.4em]
  \item Berapa byte yang dipakai bila \perintah{metode\_bayar},
        \perintah{status}, dan \perintah{jenis\_pembeli} disimpan sebagai kolom
        teks pada setiap baris fakta?
  \item Berapa byte yang dipakai bila ketiganya digantikan satu kolom
        \perintah{INTEGER}?
  \item Pada dataset ukuran sedang berisi 5 juta baris fakta, berapa selisih
        totalnya dalam megabyte?
\end{enumerate}

Perhitungan lebar baris yang lebih cermat --- termasuk header, pointer, dan
padding --- menjadi bahan Modul 5.

\section{Tugas Individu}

\subsection{Pekerjaan Wajib}
Rancang bridge table untuk hubungan banyak-ke-banyak antara promosi dan produk,
termasuk kunci, bobot jika diperlukan, dan contoh penggunaannya.

Kumpulkan \berkas{bridge-promosi-produk.md} yang memuat:

\begin{enumerate}[leftmargin=1.4em]
  \item DDL bridge table beserta kunci dan constraint-nya;
  \item penjelasan mengapa hubungan ini tidak dapat diselesaikan dengan
        menambahkan \perintah{promosi\_key} pada \perintah{fakta\_penjualan};
  \item satu contoh query yang menjawab ``berapa nilai penjualan produk yang
        tercakup promosi X'', beserta penjelasan bagaimana query itu menghindari
        penghitungan ganda;
  \item pembahasan apakah bridge table Anda memerlukan kolom bobot
        (\istilah{allocation factor}), dan apa akibatnya bagi penjumlahan
        measure bila bobot itu tidak ada.
\end{enumerate}

\begin{catatan}
Butir ketiga adalah inti tugas ini. Satu baris fakta yang produknya tercakup
tiga promosi akan terhitung tiga kali bila di-join begitu saja lewat bridge.
Ada beberapa cara menanganinya --- membatasi query pada satu promosi,
memakai \perintah{EXISTS} alih-alih \perintah{JOIN}, atau membagi measure
menurut bobot. Pilih satu, lalu jelaskan konsekuensinya.
\end{catatan}

\subsection{Pertanyaan Analisis}

\begin{enumerate}[leftmargin=1.4em]
  \item Bagian pemasaran memberi tahu bahwa harga sebuah produk sebenarnya
        sudah naik sejak dua bulan lalu, tetapi baru dilaporkan hari ini.
        Bagaimana Anda menyisipkan versi yang berlaku surut, dan apa yang harus
        dilakukan terhadap baris fakta yang sudah terlanjur menunjuk versi
        lama?
  \item Sebuah produk dihentikan penjualannya. Apakah barisnya dihapus dari
        \perintah{dim\_produk}? Jelaskan akibatnya bagi laporan historis, dan
        usulkan penanganan yang lebih baik.
  \item Anda menerapkan Type 1 pada \perintah{nama\_produk}. Setahun kemudian
        seseorang bertanya mengapa laporan lama yang sudah dicetak memuat nama
        yang berbeda dari laporan hari ini. Jelaskan sebabnya, dan sebutkan
        kapan Type 1 sebaiknya tidak dipakai meskipun nilai lamanya keliru.
\end{enumerate}

\section{Format Pengumpulan dan Checklist}

Kumpulkan DDL, skrip pemutakhiran, hasil query historis, dan rancangan bridge.
Buktikan bahwa constraint menolak versi yang tumpang tindih.

Seluruh berkas diletakkan pada \berkas{modul-03-scd2/}.

\begin{ceklis}
  \item \berkas{pre-lab.md}
  \item \berkas{kebijakan-scd.md} --- tabel kebijakan per atribut beserta alasan
  \item \berkas{ddl-scd2.sql}
  \item \berkas{load-scd2.sql}
  \item \berkas{bukti-constraint.md} --- salinan pesan galat penolakan
  \item \berkas{bukti-tiga-versi.md}
  \item \berkas{bridge-promosi-produk.md}
  \item \berkas{jawaban-analisis.md}
\end{ceklis}

\section{Troubleshooting}

\begin{table}[htbp]
\centering
\small
\caption{Masalah yang lazim muncul pada Modul 3 dan penanganannya}
\label{tab:m3-troubleshooting}
\begin{tabular}{@{}p{0.30\textwidth}>{\raggedright\arraybackslash}p{0.62\textwidth}@{}}
\toprule
\textbf{Gejala} & \textbf{Penanganan} \\
\midrule
\perintah{data type integer has no default operator class for access method
  gist} &
  Ekstensi \perintah{btree\_gist} belum aktif pada basis data ini. Jalankan
  \perintah{CREATE EXTENSION btree\_gist;} lebih dahulu. \\
\perintah{conflicting key value violates exclusion constraint} saat memasang
  constraint &
  Dimensi sudah memuat rentang yang beririsan sebelum constraint dipasang.
  Temukan barisnya dengan query celah-dan-irisan, perbaiki, lalu pasang
  kembali. \\
\perintah{duplicate key value violates unique constraint
  "uq\_dim\_produk\_kini"} &
  Langkah menutup versi lama tidak dijalankan, atau dijalankan pada baris yang
  keliru. Pastikan \perintah{UPDATE} dan \perintah{INSERT} berada dalam satu
  transaksi. \\
Versi baru tidak terbentuk padahal atribut berubah &
  Perbandingan memakai \perintah{<>} sehingga perubahan dari atau menuju
  \perintah{NULL} tidak terdeteksi. Ganti dengan
  \perintah{IS DISTINCT FROM}. \\
Jumlah baris \perintah{dim\_produk} membengkak setiap pemuatan &
  Skrip tidak idempoten: setiap eksekusi membuat versi baru meskipun tidak ada
  perubahan. Tambahkan syarat perbandingan atribut pada
  \perintah{UPDATE}. \\
Query historis mengembalikan dua baris untuk satu tanggal &
  Batas rentang ditulis tertutup di kedua ujung. Gunakan
  \perintah{>= mulai\_berlaku} dan \perintah{< selesai\_berlaku}. \\
Total penjualan berubah setelah pemuatan ulang fakta &
  Pencarian kunci memakai \perintah{baris\_kini}, bukan rentang tanggal
  transaksi. Perbaiki join pada pemuatan fakta. \\
\bottomrule
\end{tabular}
\end{table}

\section{Ringkasan}

Modul ini menjawab pertanyaan yang membedakan gudang data dari salinan basis
data operasional: \emph{apa yang berlaku pada saat itu?} Jawabannya memerlukan
tiga hal yang dibangun hari ini --- kolom rentang berlaku, constraint yang
menjaga rentang itu tetap masuk akal, dan pemuatan yang menutup versi lama alih-alih
menimpanya.

Tiga gagasan perlu dibawa ke Modul 4. \emph{Pertama}, surrogate key Modul 2 baru
sekarang terbukti gunanya: tanpa pemisahan itu, satu produk tidak mungkin
memiliki tiga baris dimensi. \emph{Kedua}, constraint bukan hiasan melainkan
alat kerja --- \perintah{EXCLUDE} menemukan kesalahan pada detik pemuatan, bukan
berbulan-bulan kemudian lewat laporan yang ganjil. \emph{Ketiga}, pencarian
kunci saat memuat fakta kini bergantung pada waktu, dan bentuk join yang keliru
tidak menimbulkan galat apa pun --- hanya laporan yang diam-diam salah.

Modul 4 memakai dimensi yang sudah tahan waktu ini untuk membangun tiga jenis
tabel fakta sekaligus, lalu menyandingkan hasilnya lewat dimensi yang dipakai
bersama.
